\makeatletter
\expandafter\gdef\csname CGTip@sec:introduction\endcsname{Section 1: Introduction (p. 2)}
\expandafter\gdef\csname CGTip@subsec:clock_intro\endcsname{Section 1.1: The Riemann--von Mangoldt clock and the ordinates (p. 2)}
\expandafter\gdef\csname CGTip@eq:intro_rvm\endcsname{Eq. (1.1), p. 2: N(T)=N\_0(T)+(1)/(pi)arg zeta ((1)/(2)+iT)+O((1)/(T)), N\_0(T)=(T)/(2pi)log(T)/(2pi e)+(7)/(8),}
\expandafter\gdef\csname CGTip@eq:intro_clock\endcsname{Eq. (1.2), p. 3: N\_0(tau\_k)=k-(1)/(2), tau\_k=(2pi (k-(11)/(8)))/(W\_0((k-(11)/(8))/e))(k>=1),}
\expandafter\gdef\csname CGTip@eq:intro_defect\endcsname{Eq. (1.3), p. 3: s\_k=k-(1)/(2)-N\_0(gamma\_k), u\_k(X)=s\_k-S\_X(gamma\_k), S\_X(t)=-(1)/(pi)sum\_p<= Xarg (1-p\textasciicircum{}-1/2-it).}
\expandafter\gdef\csname CGTip@subsec:known_intro\endcsname{Section 1.2: What is known and only used (p. 3)}
\expandafter\gdef\csname CGTip@subsec:main_results\endcsname{Section 1.3: What this paper adds (p. 3)}
\expandafter\gdef\csname CGTip@eq:intro_main\endcsname{Eq. (1.4), p. 4: (1)/(no.I\_T)sum\_kin I\_Tu\_k(X)\textasciicircum{}2=(1)/(2pi\textasciicircum{}2)log(1)/(c)+O(1+sqrt(log(1/c))).}
\expandafter\gdef\csname CGTip@sec:notation\endcsname{Section 2: Notation (p. 4)}
\expandafter\gdef\csname CGTip@subsec:symbols\endcsname{Section 2.1: List of symbols (p. 4)}
\expandafter\gdef\csname CGTip@subsec:counting\endcsname{Section 2.2: Counting functions (p. 6)}
\expandafter\gdef\csname CGTip@eq:RvM\endcsname{Eq. (2.1), p. 6: N(t)=(1)/(pi)vartheta(t)+1+S(t), S(t)=(1)/(pi)arg zeta ((1)/(2)+it)=O(log t),}
\expandafter\gdef\csname CGTip@eq:N0_def\endcsname{Eq. (2.2), p. 6: Theta\_0(t)=(t)/(2)log(t)/(2pi e)-(pi)/(8), N\_0(t)=(1)/(pi)Theta\_0(t)+1, varpi(t)=vartheta(t)-Theta\_0(t).}
\expandafter\gdef\csname CGTip@eq:N_decomposition\endcsname{Eq. (2.3), p. 6: N(t)=N\_0(t)+S(t)+(varpi(t))/(pi), L(t)=N\_0'(t)=(1)/(2pi)log(t)/(2pi).}
\expandafter\gdef\csname CGTip@eq:IT_size\endcsname{Eq. (2.4), p. 6: no.I\_T=TL(T)+O(T).}
\expandafter\gdef\csname CGTip@subsec:clock\endcsname{Section 2.3: The RvM clock and the displacements (p. 6)}
\expandafter\gdef\csname CGTip@eq:clock\endcsname{Eq. (2.5), p. 6: tau\_x=(2pi (x-(11)/(8)))/(W\_0((x-(11)/(8))/e)), (dtau\_x)/(dx)=(1)/(L(tau\_x))=(2pi)/(log(tau\_x/2pi)),}
\expandafter\gdef\csname CGTip@def:displacements\endcsname{Definition 2.1 (p. 7): For k>=1, s\_k=k-(1)/(2)-N\_0(gamma\_k) and delta\_k=tau\_k-gamma\_k.}
\expandafter\gdef\csname CGTip@eq:s_delta\endcsname{Eq. (2.6), p. 7: s\_k=L(xi\_k) delta\_kfor some xi\_k between gamma\_k and tau\_k .}
\expandafter\gdef\csname CGTip@lem:s_bounds\endcsname{Lemma 2.2 (p. 7): For every k>=1, N(gamma\_k\textasciicircum{}-)-N\_0(gamma\_k)+(1)/(2)<= s\_k<= N(gamma\_k\textasciicircum{}+)-N\_0(gamma\_k)-(1)/(2), and |s\_k|<=max(|S(gamma\_k\textasciicircum{}-)|,|S(gamma\_k\textasciicircum{}+)|)+(1)/(2)+(1)/(100).}
\expandafter\gdef\csname CGTip@eq:s_bounds\endcsname{Eq. (2.7), p. 7: N(gamma\_k\textasciicircum{}-)-N\_0(gamma\_k)+(1)/(2)<= s\_k<= N(gamma\_k\textasciicircum{}+)-N\_0(gamma\_k)-(1)/(2),}
\expandafter\gdef\csname CGTip@subsec:euler_phase\endcsname{Section 2.4: Euler phases and the corrected clock (p. 7)}
\expandafter\gdef\csname CGTip@eq:SA_def\endcsname{Eq. (2.8), p. 7: S\_A(t)=sum\_nin Aa\_nsin(tlog n), v\_A=(1)/(2)sum\_nin Aa\_n\textasciicircum{}2 ;}
\expandafter\gdef\csname CGTip@eq:SX_def\endcsname{Eq. (2.9), p. 7: S\_X(t)=-(1)/(pi)sum\_p<= Xarg (1-p\textasciicircum{}-1/2-it)=-(1)/(pi)sum\_p<= Xsum\_m>=1(sin(mtlog p))/(m p\textasciicircum{}m/2)=S\_A\_X(t).}
\expandafter\gdef\csname CGTip@eq:sigma_M_def\endcsname{Eq. (2.10), p. 7: |S\_X(t)|<=sigma\_X:=(1)/(pi)sum\_p<= Xsum\_m>=1(1)/(m p\textasciicircum{}m/2), |S\_X'(t)|<= M\_X:=(1)/(pi)sum\_p<= X(log p)/(sqrt(p)-1),}
\expandafter\gdef\csname CGTip@def:defect\endcsname{Definition 2.3 (p. 7): For X>=2 let F\_X=N\_0+S\_X and u\_k(X)=k-(1)/(2)-F\_X(gamma\_k)=s\_k-S\_X(gamma\_k). The one-step clock and its residual are tau\_k\textasciicircum{}X=tau\_k-(S\_X(tau\_k))/(L(tau\_k)), delta\_k\textasciicircum{}X=tau\_k\textasciicircum{}X-gamma\_k, r\_k(X)=L(tau\_k) delta\_k\textasciicircum{}X .}
\expandafter\gdef\csname CGTip@eq:one_step_clock\endcsname{Eq. (2.11), p. 7: tau\_k\textasciicircum{}X=tau\_k-(S\_X(tau\_k))/(L(tau\_k)), delta\_k\textasciicircum{}X=tau\_k\textasciicircum{}X-gamma\_k, r\_k(X)=L(tau\_k) delta\_k\textasciicircum{}X .}
\expandafter\gdef\csname CGTip@sec:known_results\endcsname{Section 3: External results (p. 8)}
\expandafter\gdef\csname CGTip@thm:known_basic\endcsname{Theorem 3.1 ([ref, Thms. 9.4, 9.9 and S 9.12], [ref]) (p. 8): S(t)=O(log t); gamma\_k+1-gamma\_k is bounded; and S\_1(t)=int\_0\textasciicircum{}tS(y) dy=O(log t).}
\expandafter\gdef\csname CGTip@thm:known_BW\endcsname{Theorem 3.2 (Bellotti--Wong [ref]) (p. 8): For T>= e, |N(T)-(T)/(2pi)log(T)/(2pi e)|<=BW(T).}
\expandafter\gdef\csname CGTip@thm:known_selberg_moments\endcsname{Theorem 3.3 (Selberg [ref]; see [ref, p. 225]) (p. 8): For every integer j>=1, int\_T\textasciicircum{}2TS(t)\textasciicircum{}2j dt=kappa\_j T(log log T)\textasciicircum{}j+O\_j(T(log log T)\textasciicircum{}j-1/2).}
\expandafter\gdef\csname CGTip@thm:known_korolev\endcsname{Theorem 3.4 (Selberg; Korolev [ref], [ref, Thm. 6]) (p. 8): Let Delta\_k be the integer such that g\_k+Delta\_k-1<gamma\_k<= g\_k+Delta\_k, and let 0<varepsilon\_0<10\textasciicircum{}-3. For every fixed j>=1, as n\_0->inf, uniformly for n\_0\textasciicircum{}27/82+0.9varepsilon\_0<= h\_0<= n\_0\textasciicircum{}27/82+varepsilon\_0, sum\_n\_0<k<= n\_0+h\_0Delta\_k\textasciicircum{}2j=kappa\_jh\_0(log log n\_0)\textasciicircum{}j(1+o(1)), |sum\_n\_0<k<= n\_0+h\_0Delta\_k\textasciicircum{}2j-1|<<\_j h\_0...}
\expandafter\gdef\csname CGTip@eq:korolev_moments\endcsname{Eq. (3.1), p. 8: sum\_n\_0<k<= n\_0+h\_0Delta\_k\textasciicircum{}2j=kappa\_jh\_0(log log n\_0)\textasciicircum{}j(1+o(1)), |sum\_n\_0<k<= n\_0+h\_0Delta\_k\textasciicircum{}2j-1|<<\_j h\_0(log log n\_0)\textasciicircum{}j-1.}
\expandafter\gdef\csname CGTip@thm:known_selberg_approx\endcsname{Theorem 3.5 (Selberg's approximation in the form of Fujii [ref, S 3, pp. 232--233]) (p. 8): int\_T\textasciicircum{}2T(S(t)-S\textasciicircum{}rm pr\_Y\_0(t))\textasciicircum{}4dt<< T, where Y\_0=T\textasciicircum{}1/160.}
\expandafter\gdef\csname CGTip@thm:known_MV\endcsname{Theorem 3.6 (Montgomery--Vaughan [ref, Cor. 3]) (p. 8): For complex z\_n, int\_T\textasciicircum{}2T|sum\_n<= Yz\_nn\textasciicircum{}-it|\textasciicircum{}2dt=sum\_n<= Y|z\_n|\textasciicircum{}2(T+O(n)) with an absolute implied constant.}
\expandafter\gdef\csname CGTip@thm:known_mertens\endcsname{Theorem 3.7 (Mertens [ref]) (p. 8): For y>=2, sum\_p<= y(log p)/(p)=log y+O(1) and sum\_p<= y(1)/(p)=log log y+O(1).}
\expandafter\gdef\csname CGTip@thm:known_KL\endcsname{Theorem 3.8 (Kusmin--Landau [ref, Thm. 2.1]) (p. 8): Let Omega be real with continuous monotone derivative on [y\_1,y\_2] and ||Omega'||>=1/H\_rm KL>0 there, where ||*|| is the distance to the nearest integer. Then sum\_y\_1<n<= y\_2e\textasciicircum{}2pi iOmega(n)<< H\_rm KL.}
\expandafter\gdef\csname CGTip@thm:known_landau\endcsname{Theorem 3.9 (Landau [ref]; uniform version Gonek [ref]) (p. 8): For fixed integer n>=2, sum\_0<gamma <= tn\textasciicircum{}rho=-(t)/(2pi)Lambda(n)+O(log t), zeros counted with multiplicity.}
\expandafter\gdef\csname CGTip@rem:korolev\endcsname{Remark 3.10 (Known consequences for the clock) (p. 9): (a) Moments and law of s\_k. Since vartheta(gamma\_k)/pi in(k+Delta\_k-1,k+Delta\_k] and vartheta/pi=N\_0-1+varpi/pi, one has s\_k+Delta\_k=k+Delta\_k-(1)/(2)-N\_0(gamma\_k)in[-(3)/(2)+(varpi(gamma\_k))/(pi),-(1)/(2)+(varpi(gamma\_k))/(pi)), so s\_k=-Delta\_k+ d\_k with | d\_k|<2. The index set I\_T is the integer interval (N(T),N(2...}
\expandafter\gdef\csname CGTip@eq:s_moments\endcsname{Eq. (3.2), p. 9: < s\_k\textasciicircum{}2j>\_I\_T=kappa\_j(log log T)\textasciicircum{}j(1+o(1)), < s\_k\textasciicircum{}2j-1>\_I\_T=O\_j((log log T)\textasciicircum{}j-1).}
\expandafter\gdef\csname CGTip@eq:BW_pointwise\endcsname{Eq. (3.3), p. 9: -(BW(gamma\_k)+3/8)/(L(tau\_k))<=delta\_k<=(BW(gamma\_k)-11/8)/(L(gamma\_k))(k>=1).}
\expandafter\gdef\csname CGTip@sec:identities\endcsname{Section 4: Theorem A: an exact defect identity at the zeros (p. 9)}
\expandafter\gdef\csname CGTip@thm:identity_general\endcsname{Theorem 4.1 (p. 9): Let 1<= k\_1<k\_2 be indices, F a real C\textasciicircum{}1 function on an open interval containing [gamma\_k\_1,gamma\_k\_2], and u\_k=k-(1)/(2)-F(gamma\_k) for k\_1<= k<= k\_2. For every integer j>=1, int\_gamma\_k\_1\textasciicircum{}gamma\_k\_2(N(t)-F(t))\textasciicircum{}2jF'(t) dt=sum\_k=k\_1+1\textasciicircum{}k\_2-1Psi\_j(u\_k)+(((1)/(2)+u\_k\_1)\textasciicircum{}2j+1+((1)/(2)-u\_k\_2)\textasciicircum{}2j+1)/(2j+1), where Psi\_j(u)=...}
\expandafter\gdef\csname CGTip@eq:identity_general\endcsname{Eq. (4.1), p. 9: int\_gamma\_k\_1\textasciicircum{}gamma\_k\_2(N(t)-F(t))\textasciicircum{}2jF'(t) dt=sum\_k=k\_1+1\textasciicircum{}k\_2-1Psi\_j(u\_k)+(((1)/(2)+u\_k\_1)\textasciicircum{}2j+1+((1)/(2)-u\_k\_2)\textasciicircum{}2j+1)/(2j+1),}
\expandafter\gdef\csname CGTip@eq:Psi_def\endcsname{Eq. (4.2), p. 9: Psi\_j(u)=(((1)/(2)+u)\textasciicircum{}2j+1+((1)/(2)-u)\textasciicircum{}2j+1)/(2j+1)=sum\_l=0\textasciicircum{}jbinom(2j+1,2l)(4\textasciicircum{}l-j)/(2j+1) u\textasciicircum{}2l=u\textasciicircum{}2j+sum\_l<jbinom(2j+1,2l)(4\textasciicircum{}l-j)/(2j+1)u\textasciicircum{}2l.}
\expandafter\gdef\csname CGTip@eq:identity_quadratic\endcsname{Eq. (4.3), p. 10: sum\_k=k\_1+1\textasciicircum{}k\_2-1u\_k\textasciicircum{}2=int\_gamma\_k\_1\textasciicircum{}gamma\_k\_2(N(t)-F(t))\textasciicircum{}2F'(t) dt-(k\_2-k\_1-1)/(12)-(((1)/(2)+u\_k\_1)\textasciicircum{}3+((1)/(2)-u\_k\_2)\textasciicircum{}3)/(3) .}
\expandafter\gdef\csname CGTip@eq:identity_one_gap\endcsname{Eq. (4.4), p. 10: int\_gamma\_k\textasciicircum{}gamma\_k+1(k-F(t))\textasciicircum{}2jF'(t) dt =((k-F(gamma\_k))\textasciicircum{}2j+1-(k-F(gamma\_k+1))\textasciicircum{}2j+1)/(2j+1); =(((1)/(2)+u\_k)\textasciicircum{}2j+1+((1)/(2)-u\_k+1)\textasciicircum{}2j+1)/(2j+1),}
\expandafter\gdef\csname CGTip@rem:twelfth\endcsname{Remark 4.2 (p. 10): (1)/(12) is the mean square of a centred unit sawtooth: between consecutive zeros N-F decreases through one unit in the variable F(t). If F' changes sign, (4.3) remains exact with the integral counted with sign. With F=N\_0 and Theorem 3.3, induction on j through the lower-degree terms of Psi\_j recovers the even mome...}
\expandafter\gdef\csname CGTip@cor:identity_clock\endcsname{Corollary 4.3 (p. 10): For X>=2 and 1<= k\_1<k\_2, sum\_k=k\_1+1\textasciicircum{}k\_2-1u\_k(X)\textasciicircum{}2= int\_gamma\_k\_1\textasciicircum{}gamma\_k\_2(S-S\_X+(varpi)/(pi))\textasciicircum{}2(L+S\_X')dt; -(k\_2-k\_1-1)/(12)-(((1)/(2)+u\_k\_1(X))\textasciicircum{}3+((1)/(2)-u\_k\_2(X))\textasciicircum{}3)/(3) .}
\expandafter\gdef\csname CGTip@eq:identity_LeClair\endcsname{Eq. (4.5), p. 10: sum\_k=k\_1+1\textasciicircum{}k\_2-1u\_k(X)\textasciicircum{}2= int\_gamma\_k\_1\textasciicircum{}gamma\_k\_2(S-S\_X+(varpi)/(pi))\textasciicircum{}2(L+S\_X')dt; -(k\_2-k\_1-1)/(12)-(((1)/(2)+u\_k\_1(X))\textasciicircum{}3+((1)/(2)-u\_k\_2(X))\textasciicircum{}3)/(3) .}
\expandafter\gdef\csname CGTip@lem:first_moment_identity\endcsname{Lemma 4.4 (First-moment identity) (p. 11): Let 2pi<t\_1<t\_2 be points which are not ordinates and psi in C\textasciicircum{}1[t\_1,t\_2], real or complex. Then sum\_t\_1<gamma\_k<= t\_2s\_kpsi(gamma\_k)=int\_t\_1\textasciicircum{}t\_2G(t)psi(t)L(t) dt+(1)/(2)[psi G\textasciicircum{}2]\_t\_1\textasciicircum{}t\_2-(1)/(2)int\_t\_1\textasciicircum{}t\_2G(t)\textasciicircum{}2psi'(t) dt .}
\expandafter\gdef\csname CGTip@eq:first_moment_identity\endcsname{Eq. (4.6), p. 11: sum\_t\_1<gamma\_k<= t\_2s\_kpsi(gamma\_k)=int\_t\_1\textasciicircum{}t\_2G(t)psi(t)L(t) dt+(1)/(2)[psi G\textasciicircum{}2]\_t\_1\textasciicircum{}t\_2-(1)/(2)int\_t\_1\textasciicircum{}t\_2G(t)\textasciicircum{}2psi'(t) dt .}
\expandafter\gdef\csname CGTip@eq:first_moment_measures\endcsname{Eq. (4.7), p. 11: d(G\textasciicircum{}2)=-2GL dt+sum\_t\_1<gamma<t\_22G\textasciicircum{}*(gamma)nu(gamma) delta\_gamma ,}
\expandafter\gdef\csname CGTip@sec:main_theorem\endcsname{Section 5: Theorem B: how much prime information the clock needs (p. 11)}
\expandafter\gdef\csname CGTip@thm:main\endcsname{Theorem 5.1 (p. 11): There are absolute constants T\_0 and C such that for T>= T\_0 and every X with 2<= X<= T\textasciicircum{}1/2, writing c=log X/log T, |(1)/(no.I\_T)sum\_kin I\_Tu\_k(X)\textasciicircum{}2-(1)/(2pi\textasciicircum{}2)log(1)/(c)|<= C(1+sqrt(log(1/c))).}
\expandafter\gdef\csname CGTip@eq:main_theorem\endcsname{Eq. (5.1), p. 11: |(1)/(no.I\_T)sum\_kin I\_Tu\_k(X)\textasciicircum{}2-(1)/(2pi\textasciicircum{}2)log(1)/(c)|<= C(1+sqrt(log(1/c))).}
\expandafter\gdef\csname CGTip@cor:fixed_cutoff\endcsname{Corollary 5.2 (fixed cutoff) (p. 11): For fixed X>=2, (1)/(no.I\_T)sum\_kin I\_Tu\_k(X)\textasciicircum{}2=(1)/(2pi\textasciicircum{}2)log log T+O\_X(sqrt(log log T)), and no.kin I\_T:|u\_k(X)|<= y\_0=o(no.I\_T) for every fixed y\_0>0.}
\expandafter\gdef\csname CGTip@cor:power_cutoff\endcsname{Corollary 5.3 (power cutoff) (p. 11): For fixed 0<c<=(1)/(2) and X=T\textasciicircum{}c, limsup\_T->inf(1)/(no.I\_T)sum\_kin I\_Tu\_k(T\textasciicircum{}c)\textasciicircum{}2<=(1)/(2pi\textasciicircum{}2)log(1)/(c)+C(1+sqrt(log(1/c))), and no.kin I\_T:|u\_k(T\textasciicircum{}c)|>w\_T=o(no.I\_T) whenever w\_T->inf. The same holds for X=N(T)\textasciicircum{}c with fixed 0<c<(1)/(2), since then X=T\textasciicircum{}c' with c'=c+O(log log T/log T).}
\expandafter\gdef\csname CGTip@eq:power_cutoff\endcsname{Eq. (5.2), p. 11: limsup\_T->inf(1)/(no.I\_T)sum\_kin I\_Tu\_k(T\textasciicircum{}c)\textasciicircum{}2<=(1)/(2pi\textasciicircum{}2)log(1)/(c)+C(1+sqrt(log(1/c))),}
\expandafter\gdef\csname CGTip@cor:dichotomy\endcsname{Corollary 5.4 (dichotomy) (p. 12): Let X=X(T) with 2<= X<= T\textasciicircum{}1/2. The mean square of u\_k(X) over I\_T is bounded as T->inf if and only if liminf\_T->inflog X/log T>0. For fixed X\_0>=2 and cin(0,(1)/(2)], (1)/(no.I\_T)sum\_kin I\_T(u\_k(X\_0)\textasciicircum{}2-u\_k(T\textasciicircum{}c)\textasciicircum{}2)=(1)/(2pi\textasciicircum{}2)log log T+O(sqrt(log log T)) .}
\expandafter\gdef\csname CGTip@eq:dichotomy\endcsname{Eq. (5.3), p. 12: (1)/(no.I\_T)sum\_kin I\_T(u\_k(X\_0)\textasciicircum{}2-u\_k(T\textasciicircum{}c)\textasciicircum{}2)=(1)/(2pi\textasciicircum{}2)log log T+O(sqrt(log log T)) .}
\expandafter\gdef\csname CGTip@rem:comparison_LeClair\endcsname{Remark 5.5 (p. 12): Since logno.p<= X=log X+O(log log X), the condition of Corollary 5.4 says that the logarithm of the number of primes must grow proportionally to log T. A fixed number of primes, as in the computations of [ref, Sect. VI], leaves almost all zeros violating the corrected equation by more than any fixed amount; the requir...}
\expandafter\gdef\csname CGTip@rem:localization\endcsname{Remark 5.6 (Localization of the ordinates) (p. 12): If M\_X<=(1)/(2)L(T), then F\_X'>=L-M\_X>=(1)/(2)L(T) on [T,2T], so F\_X is strictly increasing there. For kin I\_T with F\_X(T)<= k-(1)/(2)<= F\_X(2T) the equation F\_X(t)=k-(1)/(2) has exactly one root t\_k\textasciicircum{}Xin[T,2T] (intermediate value theorem), and u\_k(X)=F\_X(t\_k\textasciicircum{}X)-F\_X(gamma\_k) gives |gamma\_k-t\_k\textasciicircum{}X|<=2|u\_k(X)|/L(T) (mea...}
\expandafter\gdef\csname CGTip@subsec:main_lemmas\endcsname{Section 5.1: Lemmas on prime polynomials (p. 12)}
\expandafter\gdef\csname CGTip@eq:D_def\endcsname{Eq. (5.4), p. 12: W\_1(t)=-(1)/(pi)sum\_pin J(sin(tlog p))/(sqrt(p)), V=sum\_pin J(1)/(p),; W=S\textasciicircum{}rm pr\_Y\_0-S\_X=+- W\_1-S\_X\textasciicircum{}(2)-S\_X\textasciicircum{}(>=3),}
\expandafter\gdef\csname CGTip@lem:V_size\endcsname{Lemma 5.7 (p. 12): V=log(1/c)+O(1) uniformly for log 2/log T<= c<=(1)/(2).}
\expandafter\gdef\csname CGTip@lem:trig_MV\endcsname{Lemma 5.8 (p. 12): For real z\_n, 2<= n<= Y, int\_T\textasciicircum{}2T(sum\_n<= Yz\_nsin(tlog n))\textasciicircum{}2dt=(T)/(2)sum\_n<= Yz\_n\textasciicircum{}2+O(sum\_n<= Ynz\_n\textasciicircum{}2+(sum\_n<= Y|z\_n|)\textasciicircum{}2), and the same with cos in place of sin.}
\expandafter\gdef\csname CGTip@eq:trig_MV\endcsname{Eq. (5.5), p. 12: int\_T\textasciicircum{}2T(sum\_n<= Yz\_nsin(tlog n))\textasciicircum{}2dt=(T)/(2)sum\_n<= Yz\_n\textasciicircum{}2+O(sum\_n<= Ynz\_n\textasciicircum{}2+(sum\_n<= Y|z\_n|)\textasciicircum{}2),}
\expandafter\gdef\csname CGTip@lem:D_moments\endcsname{Lemma 5.9 (p. 13): With absolute implied constants, int\_T\textasciicircum{}2T W\textasciicircum{}2 dt =(T)/(2pi\textasciicircum{}2)V+O(T(1+sqrt(V))), ; int\_T\textasciicircum{}2T W\textasciicircum{}4 dt << T(1+V)\textasciicircum{}2, ; int\_T\textasciicircum{}2TS\_X'(t)\textasciicircum{}2 dt << T(1+log X)\textasciicircum{}2 .}
\expandafter\gdef\csname CGTip@eq:D_second\endcsname{Eq. (5.6), p. 13: int\_T\textasciicircum{}2T W\textasciicircum{}2 dt =(T)/(2pi\textasciicircum{}2)V+O(T(1+sqrt(V))),}
\expandafter\gdef\csname CGTip@eq:D_fourth\endcsname{Eq. (5.7), p. 13: int\_T\textasciicircum{}2T W\textasciicircum{}4 dt << T(1+V)\textasciicircum{}2,}
\expandafter\gdef\csname CGTip@eq:SX_derivative\endcsname{Eq. (5.8), p. 13: int\_T\textasciicircum{}2TS\_X'(t)\textasciicircum{}2 dt << T(1+log X)\textasciicircum{}2 .}
\expandafter\gdef\csname CGTip@subsec:main_proof\endcsname{Section 5.2: Proof of Theorem 5.1 (p. 13)}
\expandafter\gdef\csname CGTip@eq:E0_bounds\endcsname{Eq. (5.9), p. 13: int\_T\textasciicircum{}2TE\_0\textasciicircum{}4 dt<< T,int\_T\textasciicircum{}2TE\_0\textasciicircum{}2 dt<= T\textasciicircum{}1/2(int\_T\textasciicircum{}2TE\_0\textasciicircum{}4 dt)\textasciicircum{}1/2<< T .}
\expandafter\gdef\csname CGTip@eq:J_asymptotic\endcsname{Eq. (5.10), p. 13: E\_2=(T)/(2pi\textasciicircum{}2)V+O(T(1+sqrt(V))),}
\expandafter\gdef\csname CGTip@eq:R_fourth\endcsname{Eq. (5.11), p. 13: int\_T\textasciicircum{}2TE\_X\textasciicircum{}4 dt<< T(1+V)\textasciicircum{}2 .}
\expandafter\gdef\csname CGTip@eq:main_step1\endcsname{Eq. (5.12), p. 14: sum\_kin I\_Tu\_k(X)\textasciicircum{}2=int\_T\textasciicircum{}2TE\_X\textasciicircum{}2L dt+int\_T\textasciicircum{}2TE\_X\textasciicircum{}2S\_X' dt-(no.I\_T)/(12)+O(T\textasciicircum{}3/4).}
\expandafter\gdef\csname CGTip@eq:main_step2\endcsname{Eq. (5.13), p. 14: |int\_T\textasciicircum{}2TE\_X\textasciicircum{}2S\_X' dt|<=(int\_T\textasciicircum{}2TE\_X\textasciicircum{}4)\textasciicircum{}1/2(int\_T\textasciicircum{}2TS\_X'\textasciicircum{}2)\textasciicircum{}1/2<< T(1+V)(1+log X).}
\expandafter\gdef\csname CGTip@eq:main_step3\endcsname{Eq. (5.14), p. 14: (1)/(no.I\_T)sum\_kin I\_Tu\_k(X)\textasciicircum{}2=(V)/(2pi\textasciicircum{}2)+O(1+sqrt(V)+(1+V)/(log T)+(1+V) c),}
\expandafter\gdef\csname CGTip@sec:accounting\endcsname{Section 6: Theorem C: exact accounting for the one-step clock (p. 14)}
\expandafter\gdef\csname CGTip@eq:two_responses\endcsname{Eq. (6.1), p. 14: R\textasciicircum{}gamma\_n(I)=(2)/(a\_n)< s\_ksin(gamma\_klog n)>\_I, R\textasciicircum{}tau\_n(I)=(2)/(a\_n)< s\_ksin(tau\_klog n)>\_I,}
\expandafter\gdef\csname CGTip@eq:Cr_def\endcsname{Eq. (6.2), p. 14: Cr\_A(I)=< S\_A(tau\_k)\textasciicircum{}2>\_I-v\_A .}
\expandafter\gdef\csname CGTip@thm:accounting\endcsname{Theorem 6.1 (Finite-window accounting) (p. 14): Let I be a finite nonempty set of indices and A a set of prime powers with sum\_nin A|a\_n|<inf. item For every k and X>=2, r\_k(X)=u\_k(X)+(S\_X(gamma\_k)-S\_X(tau\_k))+chi\_k,; |chi\_k|<=(delta\_k\textasciicircum{}2)/(2pi min(gamma\_k,tau\_k)),|S\_X(gamma\_k)-S\_X(tau\_k)|<= M\_X|delta\_k|. item For every real lambda, < s\_k\textasciicircum{}2-(s\_k-lambda S\_A(tau\_k))...}
\expandafter\gdef\csname CGTip@eq:one_step_decomposition\endcsname{Eq. (6.3), p. 14: r\_k(X)=u\_k(X)+(S\_X(gamma\_k)-S\_X(tau\_k))+chi\_k,; |chi\_k|<=(delta\_k\textasciicircum{}2)/(2pi min(gamma\_k,tau\_k)),|S\_X(gamma\_k)-S\_X(tau\_k)|<= M\_X|delta\_k|.}
\expandafter\gdef\csname CGTip@eq:exact_gain\endcsname{Eq. (6.4), p. 15: < s\_k\textasciicircum{}2-(s\_k-lambda S\_A(tau\_k))\textasciicircum{}2>\_I =sum\_nin Aa\_n\textasciicircum{}2(lambda R\textasciicircum{}tau\_n(I)-(lambda\textasciicircum{}2)/(2))-lambda\textasciicircum{}2Cr\_A(I).}
\expandafter\gdef\csname CGTip@eq:full_factor_gain\endcsname{Eq. (6.5), p. 15: < s\_k\textasciicircum{}2-(s\_k-S\_X(tau\_k))\textasciicircum{}2>\_I=D\_X(I)-Cr\_X(I),; D\_X(I)=sum\_nin A\_Xa\_n\textasciicircum{}2(R\textasciicircum{}tau\_n(I)-(1)/(2))=2< s\_kS\_X(tau\_k)>\_I-v\_X,}
\expandafter\gdef\csname CGTip@rem:accounting_meaning\endcsname{Remark 6.2 (p. 15): (6.5) is an exact sign criterion for the gain of the full Euler phase at the clock points: the gain is positive exactly when the diagonal responses outweigh the cross term. No independence of the phases tau\_klog p is assumed. A threshold R\textasciicircum{}tau\_p=(1)/(2) describes a single diagonal prime contribution and does not det...}
\expandafter\gdef\csname CGTip@rem:matrix_projection\endcsname{Remark 6.3 (Optimal trigonometric correction) (p. 15): For a finite set A of frequencies let varsigma\_n(k)=sin(tau\_klog n), Q=(<varsigma\_n\_1varsigma\_n\_2>\_I)\_n\_1,n\_2in A, b=(< s\_kvarsigma\_n(k)>\_I)\_nin A and Upsilon\_lambda=sum\_nlambda\_nvarsigma\_n. Ordinary least squares gives <(s\_k-Upsilon\_lambda(k))\textasciicircum{}2>\_I=< s\_k\textasciicircum{}2>\_I-2lambda\textasciicircum{} T b+lambda\textasciicircum{} T Qlambda, minimized by lambda= Q\textasciicircum{}d...}
\expandafter\gdef\csname CGTip@rem:transport\endcsname{Remark 6.4 (Clock point against zero) (p. 15): For a prime power n, sin(gamma\_kl\_n)-sin(tau\_kl\_n)=int\_0\textasciicircum{}1(d)/(dy)sin((tau\_k-ydelta\_k)l\_n) dy, whence, exchanging the finite mean and the integral, R\textasciicircum{}gamma\_n(I)-R\textasciicircum{}tau\_n(I)=-(2l\_n)/(a\_n)int\_0\textasciicircum{}1< s\_kdelta\_kcos ((tau\_k-ydelta\_k)l\_n)>\_I dy . For fixed n and I=I\_T the right side is O(<|s\_kdelta\_k|>\_I\_T)=O(log log T/log T...}
\expandafter\gdef\csname CGTip@eq:response_transport\endcsname{Eq. (6.6), p. 15: R\textasciicircum{}gamma\_n(I)-R\textasciicircum{}tau\_n(I)=-(2l\_n)/(a\_n)int\_0\textasciicircum{}1< s\_kdelta\_kcos ((tau\_k-ydelta\_k)l\_n)>\_I dy .}
\expandafter\gdef\csname CGTip@rem:fixed_sup\endcsname{Remark 6.5 (Fixed prime sets) (p. 15): For fixed X, |delta\_k\textasciicircum{}X-delta\_k|<=sigma\_X/L(tau\_k)->0 and |u\_k(X)-s\_k|<=sigma\_X. A fixed set of primes changes neither limsup\_k|delta\_k| nor the unboundedness of the defect.}
\expandafter\gdef\csname CGTip@sec:computations\endcsname{Section 7: Computations (p. 16)}
\expandafter\gdef\csname CGTip@subsec:data\endcsname{Section 7.1: Data and methods (p. 16)}
\expandafter\gdef\csname CGTip@subsec:clock_data\endcsname{Section 7.2: The clock and the ordinates (p. 16)}
\expandafter\gdef\csname CGTip@fig:displacement\endcsname{Figure 1 (p. 16): The displacement delta\_k=tau\_k-gamma\_k of the RvM clock and the residual tau\_k\textasciicircum{}PP-gamma\_k of the one-step clock with PP=2,3,5,7 and saturation B=0.6, for k<=2000.}
\expandafter\gdef\csname CGTip@subsec:direct_defects\endcsname{Section 7.3: Direct defects (Theorem B) (p. 17)}
\expandafter\gdef\csname CGTip@tab:direct\endcsname{Table 2 (p. 17): Mean squares of the defect u\_k(X) and of the one-step residual r\_k(X) over T<gamma\_k<=2T, X=T\textasciicircum{}c.}
\expandafter\gdef\csname CGTip@fig:direct_defects\endcsname{Figure 2 (p. 18): Mean square of u\_k(X) (filled markers) and of r\_k(X) (open markers) over T<gamma\_k<=2T against log(1/c), X=T\textasciicircum{}c, for the four heights of Table 2; dashed: (1)/(2pi\textasciicircum{}2)log(1/c).}
\expandafter\gdef\csname CGTip@subsec:accounting_data\endcsname{Section 7.4: Exact accounting (Theorem C) (p. 17)}
\expandafter\gdef\csname CGTip@tab:accounting\endcsname{Table 3 (p. 18): Exact decomposition (6.5) of the one-step gain into the diagonal term D\_X and the cross term Cr\_X.}
\expandafter\gdef\csname CGTip@fig:accounting\endcsname{Figure 3 (p. 19): Gain of the one-step clock (filled markers) and cross term Cr\_X (open markers) against the cutoff X, for five windows of ordinals; dashed: v\_X.}
\expandafter\gdef\csname CGTip@subsec:one_step_cutoff\endcsname{Section 7.5: The one-step clock against the cutoff (p. 17)}
\expandafter\gdef\csname CGTip@tab:methods\endcsname{Table 4 (p. 20): Root mean square, maximum of |tau\_k\textasciicircum{}X-gamma\_k| and number of residuals >=(1)/(2) over all K ordinates, and the same quantities in the top window (10\textasciicircum{}6,K].}
\expandafter\gdef\csname CGTip@fig:rms_cutoff\endcsname{Figure 4 (p. 20): Root mean square of tau\_k\textasciicircum{}X-gamma\_k (no saturation) against the cutoff X, by window of ordinals; dots mark the minimum in each window.}
\expandafter\gdef\csname CGTip@fig:max_window\endcsname{Figure 5 (p. 21): Maximum of |tau\_k\textasciicircum{}X-gamma\_k| in each window of ordinals for the five rules of Table 4; dotted: (1)/(2).}
\expandafter\gdef\csname CGTip@sec:conditional\endcsname{Section 8: Conditional results (p. 20)}
\expandafter\gdef\csname CGTip@subsec:hypL\endcsname{Section 8.1: Fixed frequencies under Hypothesis L (p. 21)}
\expandafter\gdef\csname CGTip@hyp:L\endcsname{Hypothesis 8.1 (Hypothesis L: fixed-frequency Landau formula on the critical line) (p. 21): For every fixed integer n>=2, sum\_0<gamma <= tn\textasciicircum{}igamma=-(t)/(2pi)(Lambda(n))/(sqrt(n))+o(t), the zeros of the Riemann zeta function being counted with multiplicity.}
\expandafter\gdef\csname CGTip@prop:L_sufficient\endcsname{Proposition 8.2 (p. 21): Hypothesis L holds if sum\_0<gamma <= t|beta-(1)/(2)|=o(t).}
\expandafter\gdef\csname CGTip@lem:clock_equidistribution\endcsname{Lemma 8.3 (p. 21): For fixed real omega !=0, sum\_kin I\_Te\textasciicircum{}iomega tau\_k<<\_omega log T. Hence for fixed integers n\_1,n\_2>=2, <sin(tau\_klog n\_1)sin(tau\_klog n\_2)>\_I\_T->(1)/(2) 1\_n\_1=n\_2.}
\expandafter\gdef\csname CGTip@thm:projection\endcsname{Theorem 8.4 (Under Hypothesis L) (p. 21): Assume Hypothesis L. For every fixed integer n>=2, < s\_ksin(gamma\_klog n)>\_I\_T->(a\_n)/(2), and the same with tau\_k in place of gamma\_k; thus R\textasciicircum{}gamma\_n(I\_T),R\textasciicircum{}tau\_n(I\_T)->1 for prime powers n.}
\expandafter\gdef\csname CGTip@eq:projection_Phi\endcsname{Eq. (8.1), p. 21: H\_n(t)=[(S(y)n\textasciicircum{}iy)/(il\_n)]\_14\textasciicircum{}t-(1)/(il\_n)sum\_14<gamma\_k<= tn\textasciicircum{}igamma\_k+(1)/(il\_n)int\_14\textasciicircum{}tn\textasciicircum{}iy(L+(varpi')/(pi))dy=(i)/(l\_n)sum\_0<gamma\_k<= tn\textasciicircum{}igamma\_k+O(log t),}
\expandafter\gdef\csname CGTip@eq:projection_partial\endcsname{Eq. (8.2), p. 22: int\_T\textasciicircum{}2TS(t)n\textasciicircum{}itL(t) dt=[ H\_nL ]\_T\textasciicircum{}2T-int\_T\textasciicircum{}2T(H\_n(t))/(2pi t) dt=-i(Lambda(n))/(2pisqrt(n) l\_n)int\_T\textasciicircum{}2TL dt+o(Tlog T).}
\expandafter\gdef\csname CGTip@cor:gain\endcsname{Corollary 8.5 (Under Hypothesis L) (p. 22): For a finite set A of prime powers and lambda in R, < s\_k\textasciicircum{}2-(s\_k-lambda S\_A(tau\_k))\textasciicircum{}2>\_I\_T->(2lambda-lambda\textasciicircum{}2)v\_A, maximal at lambda=1.}
\expandafter\gdef\csname CGTip@fig:gain_fixed\endcsname{Figure 6 (p. 22): Mean-square gain < s\_k\textasciicircum{}2>-<(s\_k-S\_PP(tau\_k))\textasciicircum{}2> of two fixed prime sets PP (full local factors) in each window of ordinals, and the constants v\_PP (dashed) of Corollary 8.5.}
\expandafter\gdef\csname CGTip@subsec:gamma_response\endcsname{Section 8.2: The response at the ordinates and its Fourier reduction (p. 22)}
\expandafter\gdef\csname CGTip@eq:ATp_STp\endcsname{Eq. (8.3), p. 22: A\_T,p=-(2pisqrt(p))/(no.I\_T)int\_T\textasciicircum{}2TG(t)sin(tlog p)L(t) dt, S\_T,p=(1)/(T)int\_T\textasciicircum{}2TS(t)\textasciicircum{}2p\textasciicircum{}it dt.}
\expandafter\gdef\csname CGTip@prop:gamma_response_exact\endcsname{Proposition 8.6 (Unconditional separation of one- and two-point terms) (p. 23): R\textasciicircum{}gamma\_p(I\_T)= A\_T,p+(pisqrt(p) Tlog p)/(no.I\_T)Re S\_T,p-(pisqrt(p))/(no.I\_T)[G\textasciicircum{}2sin(tlog p)]\_T\textasciicircum{}2T; +(pisqrt(p)log p)/(no.I\_T)int\_T\textasciicircum{}2T(G\textasciicircum{}2-S\textasciicircum{}2)cos(tlog p) dt. If theta=log p/L\_T stays in a compact subset of (0,1), the last two terms are o(1) and R\textasciicircum{}gamma\_p(I\_T)= A\_T,p+2pi\textasciicircum{}2thetasqrt(p) (1+O(1/L\_T))Re S\_T,p+o(1).}
\expandafter\gdef\csname CGTip@eq:gamma_response_exact\endcsname{Eq. (8.4), p. 23: R\textasciicircum{}gamma\_p(I\_T)= A\_T,p+(pisqrt(p) Tlog p)/(no.I\_T)Re S\_T,p-(pisqrt(p))/(no.I\_T)[G\textasciicircum{}2sin(tlog p)]\_T\textasciicircum{}2T; +(pisqrt(p)log p)/(no.I\_T)int\_T\textasciicircum{}2T(G\textasciicircum{}2-S\textasciicircum{}2)cos(tlog p) dt.}
\expandafter\gdef\csname CGTip@eq:gamma_response_scaled\endcsname{Eq. (8.5), p. 23: R\textasciicircum{}gamma\_p(I\_T)= A\_T,p+2pi\textasciicircum{}2thetasqrt(p) (1+O(1/L\_T))Re S\_T,p+o(1).}
\expandafter\gdef\csname CGTip@eq:measure_def\endcsname{Eq. (8.6), p. 23: Z\_T(omega)=int\_R e\textasciicircum{}iomega t d m\_T(t), hat f\_T(omega)=int\_R f\_T(t)e\textasciicircum{}iomega tdt .}
\expandafter\gdef\csname CGTip@thm:Fourier_exact\endcsname{Theorem 8.7 (Unconditional shifted two-point representation) (p. 23): For non-ordinates T,2T, Z\_T(omega)=sum\_T<gamma<2Tnu(gamma)e\textasciicircum{}iomega gamma-int\_T\textasciicircum{}2T(L(t)+(varpi'(t))/(pi))e\textasciicircum{}iomega t dt+S(T)e\textasciicircum{}iomega T-S(2T)e\textasciicircum{}2iomega T, Z\_T(omega)=-iomega hat f\_T(omega), so that Z\_T(0)=0 and Z\_T(omega)/omega extends continuously by -ihat f\_T; and for every real omega\_0 int\_T\textasciicircum{}2TS(t)\textasciicircum{}2e\textasciicircum{}iomega\_0t dt=-(...}
\expandafter\gdef\csname CGTip@eq:measure_expansion\endcsname{Eq. (8.7), p. 23: Z\_T(omega)=sum\_T<gamma<2Tnu(gamma)e\textasciicircum{}iomega gamma-int\_T\textasciicircum{}2T(L(t)+(varpi'(t))/(pi))e\textasciicircum{}iomega t dt+S(T)e\textasciicircum{}iomega T-S(2T)e\textasciicircum{}2iomega T,}
\expandafter\gdef\csname CGTip@eq:Fourier_convolution\endcsname{Eq. (8.8), p. 23: int\_T\textasciicircum{}2TS(t)\textasciicircum{}2e\textasciicircum{}iomega\_0t dt=-(1)/(2pi)int\_R(Z\_T(omega)Z\_T(omega\_0-omega))/(omega(omega\_0-omega)) domega,}
\expandafter\gdef\csname CGTip@eq:BT_scaled\endcsname{Eq. (8.9), p. 23: B\_T(theta)=int\_R(Z\_T(alpha L\_T)Z\_T((theta-alpha)L\_T))/(alpha(theta-alpha)) dalpha,so that S\_T,p=-(B\_T(theta))/(2pi TL\_T).}
\expandafter\gdef\csname CGTip@subsec:kernel\endcsname{Section 8.3: The kernel K\_theta (p. 24)}
\expandafter\gdef\csname CGTip@eq:kernel_piecewise\endcsname{Eq. (8.10), p. 24: K\_theta(alpha)= theta, alpha <=-1,; theta-1-alpha, -1<=alpha <=theta-1,; 0, theta-1<=alpha <=1,; alpha-1, 1<=alpha <=1+theta,; theta, alpha >=1+theta,}
\expandafter\gdef\csname CGTip@eq:kernel_intersections\endcsname{Eq. (8.11), p. 24: K\_theta(alpha)=1-h(alpha)-h(theta-alpha)-h(theta)+2max (1-max(|alpha|,|theta-alpha|,theta),0),}
\expandafter\gdef\csname CGTip@prop:kernel_integral\endcsname{Proposition 8.8 (Integral and support of the kernel) (p. 24): For 0<theta<1, int\_R(K\_theta(alpha))/(alpha(theta-alpha)) dalpha=-(2Phi(theta))/(theta), the integral being absolutely convergent, and K\_theta(alpha)=0 whenever |alpha|+|theta-alpha|+theta<2.}
\expandafter\gdef\csname CGTip@eq:kernel_integral\endcsname{Eq. (8.12), p. 24: int\_R(K\_theta(alpha))/(alpha(theta-alpha)) dalpha=-(2Phi(theta))/(theta),}
\expandafter\gdef\csname CGTip@subsec:conditional_law\endcsname{Section 8.4: The moving-prime law under Hypotheses U and W (p. 24)}
\expandafter\gdef\csname CGTip@hyp:uniform_prime\endcsname{Hypothesis 8.9 (Hypothesis U: uniform one-particle input) (p. 24): For every 0<theta\_0<theta\_1<1, uniformly for primes with theta\_0<=log p/L\_T<=theta\_1, A\_T,p=1+o(1).}
\expandafter\gdef\csname CGTip@hyp:weighted_bridge\endcsname{Hypothesis 8.10 (Hypothesis W: weighted arithmetic bridge) (p. 24): Uniformly for primes with theta in any compact subset of (0,1), int\_R(Z\_T(alpha L\_T)Z\_T((theta-alpha)L\_T)+ N\_Tp\textasciicircum{}-1/2 K\_theta(alpha))/(alpha(theta-alpha)) dalpha=o((TL\_T)/(sqrt(p))), with the continuous extensions of Z\_T(omega)/omega at omega=0.}
\expandafter\gdef\csname CGTip@eq:weighted_arithmetic_bridge\endcsname{Eq. (8.13), p. 24: int\_R(Z\_T(alpha L\_T)Z\_T((theta-alpha)L\_T)+ N\_Tp\textasciicircum{}-1/2 K\_theta(alpha))/(alpha(theta-alpha)) dalpha=o((TL\_T)/(sqrt(p))),}
\expandafter\gdef\csname CGTip@prop:conditional_universal\endcsname{Proposition 8.11 (Moving-prime law under Hypotheses U and W) (p. 25): Under Hypothesis W, S\_T,p=-(1)/(2pi\textasciicircum{}2sqrt(p))(Phi(theta))/(theta)+o(p\textasciicircum{}-1/2). Under Hypotheses U and W, or U and the real part of (8.14), uniformly for theta in compact subsets of (0,1), R\textasciicircum{}gamma\_p(I\_T)=1-Phi(theta)+o(1). Conversely, under Hypothesis U, (8.15) is equivalent to the real part of (8.14). Nothing is asser...}
\expandafter\gdef\csname CGTip@eq:CTp_target\endcsname{Eq. (8.14), p. 25: S\_T,p=-(1)/(2pi\textasciicircum{}2sqrt(p))(Phi(theta))/(theta)+o(p\textasciicircum{}-1/2).}
\expandafter\gdef\csname CGTip@eq:conditional_universal\endcsname{Eq. (8.15), p. 25: R\textasciicircum{}gamma\_p(I\_T)=1-Phi(theta)+o(1).}
\expandafter\gdef\csname CGTip@rem:scale_gap\endcsname{Remark 8.12 (p. 25): Phi(theta)=2sum\_j>=1(theta\textasciicircum{}2j)/((2j-1)2j)=theta\textasciicircum{}2+(theta\textasciicircum{}4)/(6)+(theta\textasciicircum{}6)/(15)+..., so 1-theta\textasciicircum{}2 is only the quadratic truncation of (8.15), and 1-Phi(theta) vanishes at theta=0.9036.... Hypotheses U and W are not consequences of Hypothesis L. The prime polynomial alone does not give (8.14): writing S=S\textasciicircum{}rm pr\_Y+(S-S...}
\expandafter\gdef\csname CGTip@sec:heuristic\endcsname{Section 9: Heuristic support and numerical evidence (p. 25)}
\expandafter\gdef\csname CGTip@subsec:det_model\endcsname{Section 9.1: A determinantal model whose connected kernel is K\_theta (p. 25)}
\expandafter\gdef\csname CGTip@eq:real_tilt\endcsname{Eq. (9.1), p. 25: d P\_iota\_1,iota\_2=(exp -iota\_1Re Xi\_-theta-iota\_2Im Xi\_-theta)/(Eexp -iota\_1Re Xi\_-theta-iota\_2Im Xi\_-theta) d P\_0,0}
\expandafter\gdef\csname CGTip@thm:sine_kernel\endcsname{Theorem 9.1 (Connected kernel of the model) (p. 25): Let lattice frequencies tend to alpha,theta-alpha,-theta with 0<theta<1. Then R\_0\textasciicircum{}-1cum\_3(Xi\_alpha,Xi\_theta-alpha,Xi\_-theta)-> K\_theta(alpha) and R\_0\textasciicircum{}-1cum\_2(Xi\_theta,Xi\_-theta)->theta.}
\expandafter\gdef\csname CGTip@eq:third_cumulant_trace\endcsname{Eq. (9.2), p. 26: cum\_3= Tr(m\_1 m\_2 m\_3Pi\_ n)-Tr(m\_1 m\_2Pi\_ n m\_3Pi\_ n)-Tr(m\_1 m\_3Pi\_ n m\_2Pi\_ n)-Tr(m\_2 m\_3Pi\_ n m\_1Pi\_ n); +Tr(m\_1Pi\_ n m\_2Pi\_ n m\_3Pi\_ n)+Tr(m\_1Pi\_ n m\_3Pi\_ n m\_2Pi\_ n),}
\expandafter\gdef\csname CGTip@subsec:numerics_law\endcsname{Section 9.2: Numerical evidence (p. 26)}
\expandafter\gdef\csname CGTip@conj:law\endcsname{Conjecture 9.2 (p. 26): For primes with theta=log p/L\_T in compact subsets of (0,1), R\textasciicircum{}gamma\_p(I\_T)=1-Phi(theta)+o(1).}
\expandafter\gdef\csname CGTip@tab:response_errors\endcsname{Table 5 (p. 26): Errors of R\textasciicircum{}gamma\_p-(1-Phi(theta)) over the primes p<=20 000 with theta<1.}
\expandafter\gdef\csname CGTip@tab:response_bins\endcsname{Table 6 (p. 27): Bin means of R\textasciicircum{}gamma\_p and R\textasciicircum{}tau\_p in the two top windows; ``law'' is the mean of 1-Phi(theta) over the same primes, s.e. the standard error of the mean difference.}
\expandafter\gdef\csname CGTip@fig:response_law\endcsname{Figure 7 (p. 28): Bin means of R\textasciicircum{}gamma\_p (bars: standard error) against 1-Phi(theta) (solid) and 1-theta\textasciicircum{}2 (dashed), for five windows of ordinals.}
\expandafter\gdef\csname CGTip@fig:response_clock\endcsname{Figure 8 (p. 28): Bin means of the response at the clock points, R\textasciicircum{}tau\_p (squares), and at the zeros, R\textasciicircum{}gamma\_p (circles), for three windows of ordinals; solid: 1-Phi(theta).}
\expandafter\gdef\csname CGTip@sec:conclusions\endcsname{Section 10: Conclusions (p. 29)}
\expandafter\gdef\csname CGTip@q:constant\endcsname{Question 10.1 (p. 30): For fixed cin(0,(1)/(2)], do the bounded means of u\_k(T\textasciicircum{}c)\textasciicircum{}2 converge, and to (2pi\textasciicircum{}2)\textasciicircum{}-1log(1/c) plus which constant? This needs the mixed Selberg remainders and the prime-power terms, not only upper bounds.}
\expandafter\gdef\csname CGTip@q:large_c\endcsname{Question 10.2 (p. 30): Can Theorem 5.1 be extended to c>(1)/(2) (the proof uses X\textasciicircum{}2<= T in the fourth-moment step), and to a cutoff varying at every ordinate, such as X\_k=(tau\_k/2pi)\textasciicircum{}c?}
\expandafter\gdef\csname CGTip@q:prime_bridge\endcsname{Question 10.3 (p. 30): Can Hypotheses U and W be proved for moving primes, in the real or the complex form? What is the finite-height correction to 1-Phi(theta) seen in Table 6? Is there a law for the response R\textasciicircum{}tau at the clock points, with uniform control of the Gram matrix of Remark 6.3, that would predict the optimal cutoff of the com...}
\makeatother
